\subsection{Summarization and Coarsening}
\label{sec:method:reduction:summarization}
These methods compress the graph by merging nodes into super-nodes. Unlike the other two
families, which only delete, summarization must decide what a merged node \emph{is}, so
the shared rewrite that turns a clustering into a graph is described first
(\ref{sec:method:summary:mergemap}), and each method below then reduces to the single
question of which nodes it groups together.

Three methods are run, spanning domain-specific, general state of the art, and exact, so
that RQ4 has something to compare. A fourth, random merging, is a control rather than a
contender; it is what makes matched-compression comparison possible.
Table~\ref{tab:summarization_methods} lists the set.

\begin{table}[htbp]
\centering
\caption{Summarization methods.}
\label{tab:summarization_methods}
\small
\begin{tabular}{lll}
\toprule
\textbf{Method} & \textbf{Role} & \textbf{Compression} \\
\midrule
Domain-specific coarsening (TBD) & domain-specific        & method-dependent \\
Colour refinement                & adapted, exact         & fixed \\
ConvMatch                        & GNN-aware SOTA         & target ratio \\
Random within-type               & control / matching     & exact \\
\bottomrule
\end{tabular}
\end{table}

\subsubsection{Merge Map and Super-Node Schema}
\label{sec:method:summary:mergemap}
Every method produces only a \emph{clustering}: an assignment of each node to a group. One
shared rewrite turns that clustering into a graph, so that methods differ in what they merge
and in nothing else. Given a clustering, the rewrite:

\begin{itemize}
    \item sets each super-node's features to the \emph{member type counts}
    $[\#\mathrm{const}, \#\mathrm{PI}, \#\mathrm{AND}, \#\mathrm{PO}]$;
    \item remaps every edge to its endpoints' clusters, discards the resulting self-loops,
    and sums the inversion counts of parallel super-edges so that a super-edge records how
    many normal and how many inverted connections it stands for;
    \item pools the positional encoding by member \emph{minimum}, so a super-node sits at
    the level of its earliest member;
    \item records the number of discarded intra-cluster edges, and the preserved primary
    input and output counts, as graph attributes.
\end{itemize}

\paragraph{Schema compatibility.}
\label{sec:method:summary:schema}
The count representation is a strict superset of the one-hot representation: a super-node
with one member reproduces its original feature row exactly. An unreduced AIG is therefore
already a valid input in the summarized schema, without conversion. This is the property
that makes RQ5 testable with a single set of weights, and it is a design constraint on every
method rather than a convenience: a merge that produced features outside this space would
put summary and full graphs in different input spaces and make cross-state inference
undefined.

\paragraph{Boundary preservation.}
No method merges a primary input or primary output into a super-node with a different role.
The PI/PO boundary is the circuit's fixed interface, and dissolving it changes what the graph
represents rather than merely coarsening it.

% TODO: super-node content is not yet exploited; do not write as though it is.

\subsubsection{Domain-Specific Coarsening}
\label{sec:method:summary:domain}
The AIG-native slot. Two candidates are implemented; which one is run is not yet decided,
and the decision waits on measured compression and retention on the real corpus.

\begin{itemize}
    \item \textbf{Cone coarsening.} Merges along two independent axes: chains of
    single-fanout gates contract into one super-node, up to a capped chain length, reducing
    circuit \emph{depth}; and gates within a bounded band of topological levels that share
    a common immediate \emph{post}-dominator merge together, compressing parallel width
    while leaving critical-path depth untouched, so the level-based positional encoding
    stays exact. The band width and the chain cap are the only compression controls any
    domain-specific candidate offers.
    \item \textbf{MFFC skeleton coarsening.} Contracts each \emph{maximal fanout-free cone}
    into one super-node: every AND gate whose fanout lies entirely inside one cluster joins
    that cluster, iterated to a fixed point. What survives is the multi-fanout sharing
    skeleton plus the fixed PI/PO interface. It is parameter-free, so its compression is a
    property of the circuit and cannot be dialled.
\end{itemize}

The two are related rather than rival: chain contraction is the capped, chain-shaped special
case of MFFC absorption, which additionally captures cones whose interior reconverges.

MFFC coarsening invites an obvious objection, and the answer is the thesis claim in
miniature. MFFCs are prime targets for synthesis rewriting, since refactoring collapses
and re-expresses one MFFC at a time, so merging them might be expected to destroy exactly
the signal the model must read. The claim here is the opposite, and it is falsifiable:
intra-cone wiring is what local rewriting is \emph{free to change}, so it is noise with
respect to the label, whereas the sharing structure that \emph{constrains} what rewriting can
achieve is what survives the merge.

Both candidates preserve acyclicity by the same argument: a cluster is absorbed only into the
single cluster receiving all of its outgoing edges, which cannot close a cycle in a DAG\@. For
cone coarsening this holds at band width zero, where the width axis merges only within a
level; wider bands give up both the guarantee and the exactness of the level encoding.

% TODO: report that the width axis must use post-dominators, if this method is written up.
% TODO: neither cone nor MFFC compression is measured on the corpus.

\subsubsection{Graded Colour Refinement}
\label{sec:method:summary:wl}
The exact anchor of the study. Nodes are merged when they agree under $d$ rounds of colour
refinement, with $d$ set equal to the encoder's message-passing depth so that merged nodes
are exactly those the network cannot distinguish. A count-cap $c$ interpolates between two
named endpoints \citep{bollen2023exact}: at $c = \infty$ refinement distinguishes nodes by
the full multiset of neighbour colours and the merge is \emph{lossless} for a summing GNN\@;
at $c = 1$ it sees only the set of neighbour colours, which is bisimulation: coarser,
cheaper, and lossy for a counting model.

Two AIG adaptations are required. Node colours are initialised from the four-way type and
the level encoding rather than uniformly, so nodes at different depths never share a class
and the pooled positional encoding of a super-node is exact. \textbf{Edge inversion is
treated as a distinct relation} so that a normal and an inverted connection never
contribute to the same colour class. Without this, logically distinct subgraphs merge. Refinement runs backward from the primary outputs
by default; direction is a parameter, and it changes both how much compresses and which structure
survives.

At $c = \infty$ this method is orthogonal to the optimization being predicted by
construction: it removes only what the model provably cannot see, so it cannot erase the
label. That is what makes it the positive control for RQ5 rather than merely another
contender, and realising the guarantee end to end requires the schema
of~\ref{sec:method:architecture:exact}.

The obvious risk is that structural hashing has already removed the easy equivalences
(\ref{sec:prelim:algorithms:primitives}), since every corpus graph is strashed. What this
method finds is therefore \emph{residual} redundancy beyond one hop, and how much of it
exists is itself a reportable statistic about AIGs.

% TODO: run the residual-redundancy probe over depths 1 to 4; depth 1 subsumes k-SNAP \citep{tian2008ksnap}.

\subsubsection{Convolution Matching}
\label{sec:method:summary:convmatch}
The strongest general-purpose competitor: ConvMatch \citep{dickens2024convmatch} merges nodes
that are equivalent with respect to the \emph{graph convolution operation} itself rather than
with respect to a graph property. As for every other method, merges never cross node types.
The restriction is imposed on the candidate set, so the convolution objective itself is
unmodified.

It is the bar the domain-specific method must clear, and it is the right bar precisely
because it is GNN-aware but domain-blind: if an AIG-specific method beats it, the domain
claim is earned against a method optimising the same objective, not against a strawman.

Its behaviour on a deep DAG with inverted edges is untested in the literature, and one
property found here explains why it may differ from the published setting. Its merge cost is
an unweighted sum over the convolution output, so merging two low-degree nodes scores better
than merging two genuine convolution twins of higher degree, regardless of similarity.
Convolution equivalence therefore decides \emph{between} pairs of comparable degree, not
across them, which on a graph with the degree profile of an AIG biases merging toward the
input frontier.

% TODO: measure the per-graph cost of probe sampling before quoting it.

\subsubsection{Random Within-Type Merging}
\label{sec:method:summary:random}
Uniformly random nodes of the same type are merged until the node count reaches a target.
Merges are restricted to one node type so that the comparison isolates \emph{which} nodes are
merged rather than confounding it with whether the interface survives, matching the boundary
guarantee of every other method.

This serves two purposes, and the second is the important one.

As a \textbf{naive floor}, it answers the question a reader will ask first: if a
sophisticated method does not beat random merging at the same compression, its sophistication
bought nothing. It is also the exact counterpart of random edge dropout in the sparsification
family, so both halves of the study have a floor constructed the same way.

As a \textbf{matching instrument}, it is what makes RQ4 answerable at all. Compression is a
free parameter for only some methods: ConvMatch takes a target ratio directly, cone
coarsening offers only coarse integer knobs, and MFFC coarsening and colour refinement have
no compression parameter whatsoever. There is therefore no general way to make two \emph{real}
methods meet at the same ratio. Random merging reaches any target exactly by construction, so
each method is instead compared against a random control run at that method's own achieved
compression. The paired design replaces a matching problem that has no solution.

% TODO: not built. The cost is the training runs, one per matched point, not the code.
